\documentclass[11pt]{article}
\usepackage[a4paper,margin=26mm]{geometry}
\usepackage[T1]{fontenc}
\usepackage{amsmath,amssymb,amsthm,lmodern,hyperref}
\hypersetup{colorlinks=true,urlcolor=blue,linkcolor=blue}
\newtheorem{theorem}{Theorem}
\title{A polynomial mixing bound for capacity-one CYCLON}
\author{Working proof for independent review}
\date{16 November 2026}
\begin{document}
\maketitle
\begin{abstract}
For the protected/MIN mandatory-removal CYCLON model at capacity and
shuffle length one, we derive an explicit polynomial worst-case
full-state mixing bound, including arbitrary finite initial ages and
all weakly connected initial overlays. The bound counts global atomic
events. A classical logarithmic Sobolev refinement gives
$O(n^4[\log^2 n+\log(1/\varepsilon)])$; this bound remains loose.
It also bounds conditional one-ID sampling. The proof compares tree
rotations with the matroid down-up walk and uses a killed-chain lift
to handle initial directed cycles. This is an internally audited
working proof, not formal verification or external peer review. It
also gives an $\Omega(n^{3/2})$ full-state lower bound from stationary
tree geometry, leaving a substantial gap. It establishes neither an
optimal rate nor a polynomial bound for capacity at least two.
\end{abstract}

\section{Model and main result}
There are $n\ge2$ fixed labelled peers. Each cache is empty or contains
one descriptor for another peer, with a finite nonnegative integer age.
The initial directed overlay is weakly connected. At each global event,
an independent uniform peer $p$ initiates. An empty initiator does nothing.
Otherwise its singleton descriptor for $q$ is incremented and removed;
the request is the fresh descriptor $p^0$. The responder $q$ replies
with its pre-event singleton descriptor, if any, without incrementing it.
It then retains or inserts $p^0$, replacing its old descriptor as needed.
The initiator stores the reply unless it is empty or a self descriptor.
Received duplicates take minimum age. Thus the identity update is
\[
 f'(q)=p,\qquad f'(p)=
 \begin{cases}\bot,& f(q)\in\{p,\bot\},\\f(q),&\text{otherwise},\end{cases}
\]
where the right-hand values are pre-event values and all other rows stay
fixed. This is exactly the protected/MIN algorithm at $c=L=1$.
All exchanges are reliable and atomic; there is no churn or message loss.
A local query at deterministic time samples its unique cached identity,
conditional on the cache being nonempty.

\begin{theorem}
For $n\ge3$, put
\[
 N=n^{n-1},\quad H=n(n-1)^2/2,\quad G=3n^2(n-1)\min\{(n-1)\lfloor n^2/4\rfloor,\,2n-1\}.
\]
The unique stationary full-state law $\pi$ is uniform on the zero-age
rooted labelled trees. For nonnegative integers $s,t$,
\[
 \sup_x\|P^{s+t}(x,\cdot)-\pi\|_{\rm TV}
 \le\sqrt{Nn}\,e^{-s/H}
       +n(1-1/n)^s+\tfrac12\sqrt N\,e^{-t/G}.
\]
The supremum includes all weakly connected initial configurations and
arbitrary finite integer ages. For $0<\varepsilon<1$, a sufficient
mixing time is $s+t$, where
\[
 s=\left\lceil\max\left\{
 H\log(3\sqrt{Nn}/\varepsilon),\,
 n\log(3n/\varepsilon)\right\}\right\rceil,
 \quad
 t=\left\lceil G\log(3\sqrt N/(2\varepsilon))\right\rceil.
\]
This spectral bound is $O(n^5\log(n/\varepsilon))$ global events.
Section~\ref{sec:classical-lsi} improves it to
$O(n^4[\log^2 n+\log(1/\varepsilon)])$, with explicit integer times.
Conditional one-ID
error is at most $n/(n-1)$ times the full-state error whenever defined;
thus substitute $\varepsilon(n-1)/n$ above for a desired sample error
$\varepsilon$. For $n=2$, worst-case full-state distance is $2^{-t}$
for integer $t\ge1$, and available one-ID samples are already uniform.
\end{theorem}
The remainder proves the theorem. Stationary uniqueness also follows
from the uniform convergence bound, so no irreducibility claim about
the entire age-labelled state space is needed.
\section{A polynomial mixing bound at capacity one}
Let $n\ge3$, $c=L=1$, with the adopted uniform-initiator kernel.
Assume initially that the identity digraph is a rooted tree: exactly
one row is empty and every pointer leads to that root. Initial ages
are arbitrary finite nonnegative integers. The bound below does NOT
cover initial directed cycles.
Define
\[
 G=3n^2(n-1)\min\{(n-1)\lfloor n^2/4\rfloor,\,2n-1\},\quad N=n^{n-1}.
\]
The bound, in global events, is
\[
 \sup_x\|P^{s+t}(x,\cdot)-\pi\|_{\rm TV}
 \le n(1-1/n)^s+\tfrac12\sqrt N\,e^{-t/G}.
 \tag{*}
\]
In particular, for $0<\varepsilon<1$, take
\[
 s=\left\lceil n\log(2n/\varepsilon)\right\rceil,
 \qquad
 t=\left\lceil G\log(\sqrt N/\varepsilon)\right\rceil.
\]
This is $O(n^5\log(n/\varepsilon))$, deliberately unoptimized.
It is not an all-capacity polynomial bound, nor yet a polynomial
bound for arbitrary weakly connected capacity-one initial states.

\paragraph{Finite-class entrance.}
Mark all initial descriptors old and all fresh self descriptors young.
A singleton-cache activation deletes its selected partner descriptor;
a nonself reply, if present, replaces it, while the responder replaces
its singleton descriptor by the fresh initiator. Thus every old
initiator descriptor disappears, and a responder's old descriptor is
at most transferred, never duplicated. If $J$ old descriptors remain,
the next initiator holds one with probability $J/n$, so
$\mathbb E[J_{s+1}\mid X_s]\le(1-1/n)J_s$. Hence
$\Pr(J_s>0)\le n(1-1/n)^s$. Every young descriptor always has age
zero: an incremented initiator descriptor is immediately removed,
and replies are not incremented.

The rooted-tree identity class is closed under the rotations described
below. Therefore old-marker extinction implies entry into its zero-age
finite subclass. No claim is made here that the same conclusion holds
from an initial directed cycle: an off-cycle child activation can
insert a vertex into the cycle, so cycle length is not monotone.

\paragraph{The tree chain.}
Write a state as $(T,r)$, an undirected labelled spanning tree and its
root. All edges point toward $r$. Initiating $p$ with parent $q$ and
grandparent $u$ replaces $p\to q\to u$ by $q\to p\to u$.
The inverse is initiation of $q$. If $q$ is the root, the operation
merely moves the root to $p$, with the same inverse. Initiating the
root does nothing. Thus the kernel restricted to $\mathcal T$ is
symmetric, each off-diagonal edge has probability at least $1/n$,
and each state has holding probability $1/n$. Its uniform law has
$N=n^{n-1}$ states (Cayley's formula).

\paragraph{Reference chain and external input.}
Let $D$ be the uniform spanning-tree down-up chain on $K_n$: remove
a uniformly chosen tree edge and insert a uniformly chosen edge
reconnecting the resulting two components, including reinsertion
of the removed edge. The graphic matroid has rank $n-1$.
Theorem 1.1 and its matroid specialization in Anari, Liu, Oveis Gharan
and Vinzant, \emph{Log-Concave Polynomials II},
\texttt{https://arxiv.org/html/1811.01816v3}, give
$\operatorname{gap}(D)\ge1/(n-1)$. The theorem's walk is exactly
this deletion followed by conditional uniform completion.
Define $Q((T,r),(T',r'))=D(T,T')/n$; roots are refreshed uniformly.
Its stationary law is the same uniform law and its gap is at least
$1/(n-1)$ (on root-independent functions it acts as $D$, and on
root-mean-zero functions it acts as zero).

\paragraph{Simulating a reference transition.}
For a nontrivial tree exchange write $T'=T-uv+ab$, where deleting
$uv$ gives components $A$ containing $u,a$ and $B$ containing $v,b$.
Use any deterministic label convention for orientation. Perform:
\begin{enumerate}
\item Move the root from $r$ to $v$, using consecutive root-edge
reversals along the unique tree path.
\item Slide endpoint $u$ of the attachment edge $uv$ along the path
inside $A$ to $a$. At each step the next vertex $w$ is a child of
the current attachment $x$, whose parent is $v$. Initiating $w$
replaces $xv$ by $wv$ and preserves the undirected edge $wx$.
\item Move the root from $v$ to $a$ across the current attachment
edge $av$.
\item Slide endpoint $v$ along the path inside $B$ to $b$, now
with root $a$, by the same operation.
\item Move the root to $r'$ along the final tree.
\end{enumerate}
The length is at most
$2(n-1)+(|A|-1)+(|B|-1)+1=3(n-1)$.
Every intermediate unrooted tree is $T$, or $T-uv+xy$: it differs
from $T$ by at most one basis exchange. For $T=T'$, just move the
root along its tree path, which has the same properties.
Chronologically erase loops from each path, making it simple;
length and the one-exchange property are preserved. These paths
also prove irreducibility.

\paragraph{Polynomial weighted congestion.}
Decompose each $Q$ transition into its elementary down-up choices,
including deletion and reinsertion of the same edge. If deleting an
edge yields a forest $F$ with component sizes $a,b$, the elementary
transition weight, including the final root choice, is
$1/[n(n-1)ab]$. For different-tree transitions the removed edge is
unique. For same-tree transitions use the direct root-moving path
for every such elementary choice. This splitting sums to exactly $Q$.

Fix an oriented nontrivial tree-chain edge $z\to w$ and let $S$ be
the unrooted tree at $z$. Any reference path through it keeps the
same forest $F=T-e=T'-f$ throughout; in particular $S=F+g$ for some
edge $g\in S$. For this $g$, the forest $F=S-g$ is determined.
Write $a_g,b_g$ for its component sizes. There are at most
$(a_gb_g)^2$ ordered choices of reconnecting edges $(e,f)$ for the
endpoint trees, and at most $n^2$ choices of starting/finishing roots.
Each elementary path has length at most $3(n-1)$, and after loop
erasure traverses the fixed oriented edge at most once. Its stationary
starting mass is $1/N$. Since $\pi(z)P(z,w)\ge1/(Nn)$, the congestion
is at most
\[
 \begin{split}
 A&\le Nn\sum_{g\in S}
 n^2(a_gb_g)^2\frac{1}{Nn(n-1)a_gb_g}\,3(n-1)\\
 &=3n^2\sum_{g\in S}a_gb_g
 \le3n^2(n-1)\lfloor n^2/4\rfloor.
 \end{split}
\]
This count includes $e=f$, so diagonal root-refresh transitions require
no additional estimate. Paths of zero length contribute nothing.
Cauchy--Schwarz along the paths, summed over the elementary transition
decomposition, gives $\mathcal E_Q(f,f)\le A\mathcal E_P(f,f)$.
Combining with $\operatorname{gap}(Q)\ge1/(n-1)$ first gives
the denominator $3n^2(n-1)^2\lfloor n^2/4\rfloor$. The following
root-aware refinement improves it to the stated $G$.

\paragraph{Shared counting conventions.}
For a forest $F$ with two components of sizes $a,b$, each elementary
down-up move has weight $1/[n(n-1)ab]$, including the final root choice.
There are at most $n^2$ endpoint root pairs. A canonical path has length
at most $3(n-1)$ and stationary starting mass $1/N$.
The target edge capacity is at least $1/(Nn)$. Consequently, a class
with at most $M$ ordered reconnecting-edge pairs contributes at most
\[
 (Nn)\frac{1}{N}\frac{n^2 M}{n(n-1)ab}\,3(n-1)
 =\frac{3n^2 M}{ab}
\]
to congestion. Paths are chronologically loop-erased, so a fixed
oriented edge is used at most once. Erasure only removes transitions;
it cannot introduce a new stage or destroy a restriction valid for
all original occurrences of a retained transition. Summing categories
may overcount and is therefore safe.

\paragraph{Shape-changing target transition.}
Fix an oriented tree rotation $z\to w$ whose underlying trees differ.
It can occur only in one of the two attachment-slide stages, never in
a root-moving stage. During either slide, the unchanged endpoint of
the moving attachment is the current root $r$. Thus any shape-changing
target transition whose grandparent is not its root is unused.
The common forest is uniquely $F=S_z\cap S_w$. Let $C$ be its
component containing the two sliding endpoints, of size $a$, and let
the component containing the root have size $b$.
\begin{itemize}
\item In the first slide stage, the original removed edge is $\{u,r\}$
with $u\in C$, so there are at most $a$ choices for that edge.
The final reconnecting edge has at most $ab$ choices. Therefore
$M\le a^2b$, and this stage contributes at most $3n^2a$.
\item In the second slide stage, the final inserted edge has one
endpoint equal to $r$ and its other endpoint in $C$, so it has at
most $a$ choices. The original reconnecting edge has at most $ab$
choices. Again $M\le a^2b$ and the contribution is at most $3n^2a$.
\end{itemize}
The total shape-changing load is at most $6n^2a\le6n^2(n-1)$.
This count does not assume that every listed edge pair's path uses
the target transition; including additional pairs only increases it.

\paragraph{Root-reversal target transition.}
Now fix a transition with the same underlying tree $S$ and adjacent
roots $x,y$. There are three possible path stages.
\begin{enumerate}
\item During the initial root move the starting unrooted tree is
$T=S$. For each removed edge $e\in S$, its cut sizes $a,b$ allow
at most $ab$ final reconnecting edges. Thus $M\le ab$, yielding
at most $3n^2$ for each removed edge and $3n^2(n-1)$ in total.
Direct root-only paths for $e=f$ are included in this initial-stage
count; their actual lengths are smaller than the bound used here.
\item During the final root move the finishing unrooted tree is
$T'=S$. Sum over its final inserted edge $f\in S$. For each such
edge the original reconnecting edge has at most $ab$ choices, giving
another bound $3n^2(n-1)$.
\item During the central crossing the old attachment endpoint is
$v=x$ and the new endpoint is $a_{\rm path}=y$. Hence
$F=S-\{x,y\}$ is fixed. Write $a,b$ for the sizes of the components
containing $y,x$, respectively. The original removed edge has form
$\{u,x\}$ with $u$ in the first component, giving at most $a$
choices; the final inserted edge has form $\{y,b_{\rm path}\}$
with $b_{\rm path}$ in the second, giving at most $b$ choices.
Thus $M\le ab$, and this stage contributes at most $3n^2$.
\end{enumerate}
The total root-reversal load is at most
$6n^2(n-1)+3n^2=3n^2(2n-1)$, which also bounds the shape-changing
load. Multiply by the reference inverse gap $n-1$, and take the
minimum with the original bound, to obtain the stated $G$. All subsequent
negative-spectrum, killed-lift, age-extinction and conditional-sampling
arguments remain valid with this replacement.

Because $P=I/n+(1-1/n)R$ for a Markov kernel $R$, every eigenvalue
of $P$ is at least $-1+2/n$. Since $1/G\le2/n$, every nonconstant
eigenvalue has absolute value at most $1-1/G$. The elementary
finite reversible-chain $L^2$ bound now gives
\[
 \sup_{y\in\mathcal T}\|P^t(y,\cdot)-\pi\|_{\rm TV}
 \le\tfrac12\sqrt{N-1}(1-1/G)^t
 \le\tfrac12\sqrt N e^{-t/G}.
\]
Split at deterministic time $s$ according to whether all old
markers have disappeared. The bad part costs at most its probability;
the good part is a mixture of states in $\mathcal T$. This proves
$(*)$ without a stopping-time conditioning issue.

\paragraph{Sample consequence and scope.}
At stationarity exactly one row is empty, so symmetry gives
$a_1=(n-1)/n$. The sharp conditioning inequality yields one-ID
conditional error at most $n/(n-1)$ times the full-state bound,
whenever availability is positive. No corresponding polynomial
bound for $c\ge2$ follows from this argument: the rooted-tree
representation and its reversible uniform law are special to $c=1$.
\section{Polynomial entrance from a directed cycle at capacity one}
For $n\ge3$, let $N=n^{n-1}$ be the number of rooted labelled trees.
The entrance estimate below does not depend on their mixing gap. Let $\mathcal U$ be the finite
set of weakly connected outdegree-one identity configurations; these
have one directed cycle and no empty row. Let $K$ denote the actual
capacity-one identity transition matrix restricted to $\mathcal U$,
killing the chain when an empty row first appears. Every nonkilling
rotation has its inverse obtained by activating the former responder,
so $K$ is symmetric and substochastic. Ages do not affect this identity
kernel at capacity one.

\paragraph{Lift.}
Use pairs $(y,a)$ where $y=(T,r)$ is a rooted labelled tree and
$a\in V$ is a mark. Define $\Phi(y,a)$ to be the rooted tree itself
if $r=a$; otherwise add the pointer $r\to a$, producing an element
of $\mathcal U$. A unicycle with cycle length $k$ has exactly $k$
preimages, obtained by deleting any one of its directed cycle edges.
A rooted tree has exactly one preimage, with its root as mark.
In particular $|\mathcal U|\le Nn$.

Use unnormalized counting measures throughout. The killed Dirichlet form is
\[
 \mathcal E_K(f)=\frac12\sum_{u,v\in\mathcal U}
 K(u,v)(f(u)-f(v))^2
 +\sum_{u\in\mathcal U}\left(1-\sum_vK(u,v)\right)f(u)^2.
\]
Put $H=n(n-1)^2/2$.

\paragraph{Root-only lifted walk.}
For each fixed unrooted labelled tree $T$ and mark $a$, retain only
root-reversal transitions of the lifted tree chain. Each neighbor
of the current root becomes the new root with probability $1/n$;
all other probability is assigned to a self-loop. Thus the fixed
$(T,a)$ block is a symmetric nearest-neighbor walk on $T$. Its
boundary is the single root value $r=a$.
For $f$ on transient unicycles, extend $f$ by zero on rooted trees
and put $h(T,r,a)=f(\Phi(T,r,a))$.

\paragraph{Unique projected edge.}
A nonboundary lifted root reversal from root $q$ to its child $p$
projects to the actual activation $p\to q$: the added pointer
$q\to a$ is the reply and the new pointers are $q\to p,p\to a$.
If $p=a$, the self reply is discarded and this is absorption into
a rooted tree. Conversely, for an actual transient transition
initiated by $p$ with responder $q$, such a lifted edge can exist
only with lifted source root $q$ and mark equal to the old pointer
$q$'s target. These are forced. If $q$ is not on the cycle, deleting
its pointer does not yield a rooted tree and there is no preimage;
if $q$ is on the cycle, the preimage is unique.
A boundary root reversal from root $a$ to its child $p$ projects
to the reciprocal cycle $a\leftrightarrow p$; activation of $a$
in that cycle returns to the original boundary state. Thus it is
the reverse of the unique killing edge, with the same probability
$1/n$. Reversing orientation does not produce a second undirected
lifted edge. Every contributing root-only lifted edge therefore maps
injectively to an actual transient or killing edge, preserving weight.

\paragraph{Dirichlet estimate with counting measures.}
Let $\mathcal E_{\rm root}$ be the sum of raw Dirichlet forms over
all fixed $(T,a)$ blocks. The injective edge mapping gives
\[
 \mathcal E_{\rm root}(h)\le\mathcal E_K(f).
\]
For fixed $T,a$, $h(T,a,a)=0$. Along the path from any root $r$
to $a$, Cauchy--Schwarz yields
\[
 h(T,r,a)^2\le(n-1)\sum_{e\in\operatorname{path}_T(r,a)}(\Delta_eh)^2.
\]
Summing over $r$, any edge occurs in at most $n-1$ such paths, so
\[
 \sum_r h(T,r,a)^2
 \le(n-1)^2\sum_{e\in T}(\Delta_eh)^2
 =n(n-1)^2\mathcal E_{\rm root}^{T,a}(h).
\]
Summing over $T,a$ and using the lift multiplicity $k(u)\ge2$
for every transient unicycle $u$ gives
\[
 \mathcal E_K(f)\ge\mathcal E_{\rm root}(h)
 \ge\frac1{2H}\sum_{T,r,a}h(T,r,a)^2
 =\frac1{2H}\sum_u k(u)f(u)^2
 \ge\frac1H\sum_u f(u)^2.
\]
Thus the largest eigenvalue of symmetric substochastic $K$ is at
most $1-1/H$. Perron--Frobenius bounds its spectral radius by the
same number, and $|\mathcal U|\le Nn$ gives the stated survival
bound by Cauchy--Schwarz:
\[
 \sup_x\Pr_x(\tau>s)\le\sqrt{Nn}\,e^{-s/H}.
\]
This uses the factor two from the minimum cycle length, rather than discarding it.

\paragraph{Adding ages and mixing.}
The initial-old-marker expectation bound
$\Pr(J_s>0)\le n(1-1/n)^s$ is valid from both trees and unicycles;
it never relied on cycle monotonicity. All young descriptors have
age zero. Therefore at deterministic time $s$ the probability of
being outside the zero-age rooted-tree class is at most
\[
 \sqrt{Nn}\,e^{-s/H}+n(1-1/n)^s.
\]
After $t$ further events, the tree-chain spectral estimate gives
\[
 \sup_x\|P^{s+t}(x,\cdot)-\pi\|_{\rm TV}
 \le\sqrt{Nn}\,e^{-s/H}
       +n(1-1/n)^s+\tfrac12\sqrt N\,e^{-t/G}.
\]
The refined constants give a bound of $O(n^5\log(n/\varepsilon))$
for every weakly connected capacity-one initial state with finite ages.
For example choose
\[
 s=\left\lceil\max\left\{
 H\log(3\sqrt{Nn}/\varepsilon),\,
 n\log(3n/\varepsilon)\right\}\right\rceil,
 \qquad
 t=\left\lceil G\log(3\sqrt N/(2\varepsilon))\right\rceil.
\]
No assertion for capacity at least two follows from this lift.

\paragraph{Two peers.}
For $n=2,c=L=1$, an initial reciprocal pair enters the zero-age
rooted-tree class in one event, with uniformly distributed root.
From a one-edge state its initial descriptor survives only while
its holder has never initiated, with probability $2^{-t}$. Once
that descriptor disappears the two zero-age rooted states evolve
by uniform root refresh. Directly, the worst-case full-state TV
distance is $2^{-t}$ for integer $t\ge1$ (and is one at $t=0$).
Conditional one-ID output is the sole other peer and is exactly
uniform whenever it is available.

\paragraph{Conditional-sample estimate.}
Append sample extraction to a state law, using a failure symbol for an
empty cache; this cannot increase total variation. For two laws with
availability probabilities $p,q>0$ and conditional output laws $u,v$,
\[
 \|u-v\|_{\rm TV}\le\frac{\|\mu-\nu\|_{\rm TV}}{\max(p,q)}.
\]
Indeed, if $p\ge q$, overlap outside the available-output alphabet is
at most $1-p$, while overlap within it is at most
$p\sum_z\min(u_z,v_z)$. Subtracting total overlap from one proves the
inequality; the other case is symmetric. Stationary availability is
$(n-1)/n$ by uniform rooted-tree symmetry, and the conditional parent
identity is uniform over other peers by relabelling symmetry. Applying
the inequality proves the sample estimate. If full-state error is
strictly below stationary availability, transient availability is
positive automatically.

\section{Classical logarithmic Sobolev refinement}
\label{sec:classical-lsi}
Assume $n\ge3$ and retain the already proved capacity-one constants
\[
 h=\lfloor n^2/4\rfloor,\quad N=n^{n-1},\quad H=n(n-1)^2/2,
 \quad A=3n^2\min\{(n-1)h,2n-1\},\quad G=(n-1)A.
\]
Put
\[
 C=(n-1)\bigl(3+\log(nh)\bigr),\qquad \Lambda=AC.
\]
The new argument establishes a classical logarithmic Sobolev
inequality $\operatorname{Ent}_\pi(f^2)\le\Lambda\mathcal E_P(f,f)$
on zero-age rooted trees. Unlike a direct modified-log-Sobolev
comparison, this inequality uses the squared-difference Dirichlet
form for which the established canonical-path comparison is valid.

\paragraph{Elementary uniform-space inequality.}
On the uniform distribution $\nu$ on $m$ points,
\[
 \operatorname{Ent}_\nu(f^2)\le(\log m+3)\operatorname{Var}_\nu(f).
\]
To prove this, replace $f$ by $|f|$; entropy is unchanged and variance
cannot increase. Scale so $\mathbb E f^2=1$, unless $f=0$ identically.
Then $0\le f\le\sqrt m$. For $\varphi(x)=x^2\log x^2$, extended
continuously by zero at zero, $\varphi''(x)=2\log x^2+6\le2\log m+6$.
Taylor's upper bound about $\bar f=\mathbb E f\le1$, followed by
averaging, gives
\[
 \mathbb E\varphi(f)\le\varphi(\bar f)+(\log m+3)\operatorname{Var}(f)
 \le(\log m+3)\operatorname{Var}(f).
\]
The upper Taylor bound remains valid at zero by continuity, or by
integrating the upper bound on the second derivative. This proves
the inequality for all real $f$ by homogeneity.

\paragraph{Entropy factorization for the reference chain.}
Let $Q$ be the existing reference chain: a uniform spanning-tree
down-up step followed by an independent uniform root refresh. It
is a heat-bath chain with latent variable the two-component forest
$F$ obtained by deleting a uniformly chosen tree edge. Given $F$
with component sizes $a,b$, the conditional state is uniform among
$nab$ completions (an edge across the cut and a root).

We use the half-step entropy contraction in Cryan, Guo and Mousa,
\emph{Modified log-Sobolev inequalities for strongly log-concave
distributions}, Lemma 11 and equation (32),
\url{https://arxiv.org/html/1903.06081v3}. For rank $r=n-1$, the
entropy of a nonnegative function projected to the deleted-edge
forest is at most $(1-1/r)$ of its original entropy.
To apply it here, first average $f^2(T,r)$ over the root, obtaining
$g(T)$. Root averaging cannot increase entropy; then the cited
half-step contraction applies to $g$. Thus
\[
 \operatorname{Ent}_{F}\bigl(\mathbb E[f^2\mid F]\bigr)
 \le(1-1/(n-1))\operatorname{Ent}_\pi(f^2).
\]
The entropy chain rule and the uniform-space bound give
\[
 \begin{split}
 \operatorname{Ent}_\pi(f^2)
 &\le(n-1)\mathbb E_F\operatorname{Ent}(f^2\mid F)\\
 &\le(n-1)(3+\log(nh))\mathbb E_F\operatorname{Var}(f\mid F)
 =C\mathcal E_Q(f,f).
 \end{split}
\]
The last equality holds because $Q$ samples two conditionally
independent completions of the same forest in stationarity.
The previously proved comparison $\mathcal E_Q\le A\mathcal E_P$
therefore establishes the claimed classical inequality with $\Lambda=AC$.
No modified-log-Sobolev path inequality is assumed.

\paragraph{Continuous-time smoothing only as an analytic tool.}
Let $S_u=\exp(u(P-I))$. In the convention
$\alpha=\inf\mathcal E(f,f)/\operatorname{Ent}(f^2)$, the classical
hypercontractivity theorem gives
$\|S_u\|_{2\to q}\le1$ for $q-1\le e^{4\alpha u}$.
See Diaconis and Saloff-Coste, \emph{Logarithmic Sobolev inequalities
for finite Markov chains}, Theorem 3.5(ii),
\url{https://pi.math.cornell.edu/~lsc/papers/log-aoap.pdf}.
Self-adjoint duality gives $\|S_u\|_{p\to2}\le1$ for
$p=1+e^{-4u/\Lambda}$, using $\alpha\ge1/\Lambda$.
For a point-mass initial density $f_x=N\mathbf1_{\{x\}}$, choose
\[
 u_0=\frac{\Lambda}4\log\log N.
\]
Then $p-1=1/\log N$, and
$\|f_x\|_p=N^{1-1/p}\le e$. Thus
$\|S_{u_0}f_x-1\|_2\le e$. The spectral gap bound $1/G$ gives
\[
 \|S_u f_x-1\|_2\le e\,e^{-(u-u_0)/G}\qquad(u\ge u_0).
\]
These are analytic estimates; no conversion of clock units is being
used for the final discrete-event result.

\paragraph{Returning rigorously to discrete events.}
Expand $f_x-1$ in an orthonormal eigenbasis of $P$. On nonnegative
nonconstant eigenvalues $0\le\lambda<1$, the inequality
$\lambda^t\le e^{-t(1-\lambda)}$ shows that the discrete-time
positive-spectrum $L^2$ norm is at most the continuous-time norm
at $u=t$. The negative spectrum is bounded separately: the holding
probability $1/n$ implies $|\lambda|\le1-2/n$ when $\lambda<0$.
Since $\|f_x-1\|_2=\sqrt{N-1}$, its contribution is at most
$\sqrt N e^{-2t/n}$. Therefore, for integer $t\ge u_0$,
\[
 \sup_{x\in\mathcal T}\|P^t(x,\cdot)-\pi\|_{\rm TV}
 \le\tfrac12\left(e\,e^{-(t-u_0)/G}+\sqrt N e^{-2t/n}\right).
\]
For $0<\delta<1$, a sufficient rooted-class time is
\[
 T(\delta)=\left\lceil\max\left\{
 \frac{\Lambda}4\log\log N+G\log(e/\delta),\quad
 \frac n2\log(\sqrt N/\delta)\right\}\right\rceil.
\]
Each summand inside the TV upper bound is then at most $\delta$,
so their half-sum is at most $\delta$.

\paragraph{Arbitrary finite-age initial states.}
For $0<\varepsilon<1$ use the already proved entrance and old-token
bounds with
\[
 s=\left\lceil\max\left\{
 H\log(4\sqrt{Nn}/\varepsilon),\quad
 n\log(4n/\varepsilon)\right\}\right\rceil.
\]
Each of the two failure terms is at most $\varepsilon/4$.
After another $T(\varepsilon/2)$ events, the full-state TV error is
therefore at most $\varepsilon$. Since $A=O(n^3)$, $\Lambda=O(n^4\log n)$,
$G=O(n^4)$, and $H=O(n^3)$, this yields
\[
 t_{\rm mix}(\varepsilon)
 =O\!\left(n^4\left[\log^2 n+\log(1/\varepsilon)\right]\right)
\]
in global atomic events. For particular parameters one may take the
smaller of this sufficient time and the earlier spectral sufficient time.
Conditional one-ID sampling uses the same
$n/(n-1)$ factor as before. The two-peer exact result and the
$\Omega(n^{3/2})$ full-state lower bound are unchanged. No result for
capacity at least two or a matching local-query lower bound follows.

\section{A depth-based capacity-one lower bound}
Work in the zero-age rooted-tree class at $c=L=1$ with uniform
stationary law. Fix a labelled vertex $v$ and let $D$ be its tree
distance to the root. Put $m_n=\mathbb E_\pi D$ and
$v_n=\operatorname{Var}_\pi D$.
The ordinary spectral gap satisfies the exact variational bound
\[
 \operatorname{gap}(P)\le g_n:=\frac{m_n}{n v_n}.
\]
Indeed, if the current path from $v$ to the root has $D=k$ edges,
exactly the $k$ nonroot vertices on this path decrease $D$ by one
when they initiate. An off-path child of any path vertex increases
$D$ by one when it initiates. All other initiations leave $D$
unchanged. No move changes depth by more than one. If $B$ is the
number of these off-path children, conditional drift is $(B-D)/n$.
Stationarity implies $\mathbb E B=\mathbb E D$, whence
\[
 \mathcal E_P(D,D)=\tfrac12\mathbb E[(D_1-D_0)^2]
 =\frac{\mathbb E(B+D)}{2n}=\frac{m_n}{n}.
\]
Apply the variational formula for the reversible spectral gap.

For $0\le k\le n-1$, exact Cayley counting gives
\[
 \Pr_\pi(D=k)=\frac{(k+1)(n-1)_k}{n^{k+1}},
 \qquad \Pr_\pi(D\ge k)=\frac{(n-1)_k}{n^k}.
\]
To count the first expression, choose the ordered $k$ distinct
vertices after $v$ on the root path; there are $(n-1)_k$ choices.
For each prescribed path, the number of spanning trees containing
its $k$ edges is $(k+1)n^{n-k-2}$. To see this, contract the
prescribed path to a component of size $k+1$ and leave all other
vertices as singleton components. A tree on these components has
$\prod_C |C|^{\deg(C)}$ choices of original edge endpoints.
The weighted Pr\"ufer-code identity sums these weights to
$(\prod_C|C|)(\sum_C|C|)^{m-2}$, where $m$ is the number of
components, giving the stated count; the single-component case
has exactly one completion. Divide by the $n^{n-1}$ rooted trees. The tail
identity follows by telescoping, including $k=0$.
Consequently
\[
 m_n=\sum_{k=1}^{n-1}\frac{(n-1)_k}{n^k},\qquad
 \mathbb E D^2=\sum_{k=1}^{n-1}(2k-1)\frac{(n-1)_k}{n^k}.
\]
For fixed $x\ge0$, putting $k=\lfloor x\sqrt n\rfloor$ gives
$\Pr(D\ge k)\to e^{-x^2/2}$, by taking logarithms of the product.
The bound $\Pr(D\ge k)\le e^{-k(k+1)/(2n)}$ gives a summable
Gaussian majorant for the scaled Riemann sums. Hence
\[
 m_n\sim\sqrt{\pi n/2},\qquad
 v_n\sim(2-\pi/2)n,
\]
and
\[
 \operatorname{gap}(P)\le
 \left(\frac{\sqrt{2\pi}}{4-\pi}+o(1)\right)n^{-3/2}.
\]
In particular the relaxation time is $\Omega(n^{3/2})$ global events.
For $g_n<1$ and $0<\varepsilon<1/2$, the usual eigenfunction
lower bound gives
\[
 t_{\rm mix}(\varepsilon)\ge
 \frac{\log(1/(2\varepsilon))}{-\log(1-g_n)}.
\]
For completeness, choose a second-eigenvalue eigenfunction and a
state where its absolute value is maximal. Its time-$t$ expectation
differs from stationarity by $\lambda_2^t$ times that maximum, so
TV distance is at least $\lambda_2^t/2$; and
$\lambda_2\ge1-g_n$. This establishes a fixed-error full-state
mixing lower bound $\Omega(n^{3/2})$. It does not establish a
matching lower bound for a single conditional parent-identity query.

\section{Verification and limitations}
The companion standard-library checker
\begin{center}\small\path{verification/verify_c1_polynomial_comparison.py}\end{center}
enumerates rooted trees by Pr\"ufer codes and checks all elementary
basis-exchange/root pairs through five peers, together with the lift
edge mapping and its multiplicity. It checks legality, endpoints,
length and the shared-forest invariant. These are finite corroborations,
not replacements for the arbitrary-$n$ arguments above. The power nine
is a coarse consequence of the comparison; no lower bound matching it
is claimed. This result does not settle polynomial mixing when $c\ge2$.
\end{document}
